% Source: https://artofproblemsolving.com/community/c3h2380519_hopefully_theres_nothing_wrong_with_this
% Stable identity: aops-topic-2380519-q06
Let $N$ be the midpoint of $CA$ and let $B'$ be the reflection of $B$ in the point $N$. The diagonals $BB'$ and $AC$ of $ABCB'$ bisect each other, so $ABCB'$ is a parallelogram with sides $8$ and $9$ and diagonals $AC=10$ and $BB'=2\,BN$.

\begin{asy}
usepackage("amsmath");
size(5cm);
defaultpen(fontsize(10pt));
pair A=(0,0), C=(10,0);
real x=(64-81+100)/20; pair B=(x,sqrt(64-x^2));
pair Nm=(A+C)/2, B2=2*Nm-B;
draw(A--B--C--cycle,linewidth(0.9)); draw(B--B2,blue); draw(A--B2--C,dashed);
dot(Nm);
label("$A$",A,W); label("$B$",B,N); label("$C$",C,E); label("$N$",Nm,SE); label("$B'$",B2,S);
label("$8$",(A+B)/2,NW); label("$9$",(B+C)/2,NE); label("$5$",(A+Nm)/2,S); label("$5$",(Nm+C)/2,S);
\end{asy}

In a parallelogram the sum of the squares of the diagonals equals the sum of the squares of the four sides (by the law of cosines in two adjacent triangles, whose angles at a common side are supplementary). Hence
\[(2BN)^2+10^2=2\cdot8^2+2\cdot9^2=290,\qquad 4BN^2=190,\qquad BN=\dfrac{\sqrt{190}}{2}.\]
Since $190=2\cdot5\cdot19$ is squarefree, $a=190$, $b=2$ and $a+b=\boxed{192}$.
